\documentclass[11pt]{article}

% EMNLP 2026 Industry Track — double-blind review version
\usepackage[review]{acl}

\usepackage{times}
\usepackage{latexsym}
\usepackage[T1]{fontenc}
\usepackage[utf8]{inputenc}
\usepackage{microtype}
\usepackage{inconsolata}
\usepackage{graphicx}
\usepackage{booktabs}
\usepackage{amsfonts}
\usepackage{nicefrac}
\usepackage{xcolor}
\usepackage{enumitem}
\usepackage{multirow}
\usepackage{hyperref}
\usepackage{url}

\title{AgentTelemetry in Production:\\Deployment Experience with Span-Kind Observability\\Across 13 LLMs and Three Multi-Agent Topologies}

\author{Anonymous \\
  Anonymous Affiliation \\
  \texttt{anon@anon.anon} \\}

\begin{document}

\maketitle

\begin{abstract}
Production deployments of LLM-based autonomous agents fail in ways
that standard observability cannot diagnose. Building on the
\textsc{AgentTelemetry} taxonomy of nine agent-specific
OpenTelemetry span kinds~\citep{aiware2026}, we report deployment
experience from stress-testing the same instrumentation across (a) 13
production-tier LLMs from Anthropic and OpenAI on a research-assistant
agent task, (b) a Claude Opus matched-pair replication of the prior
single-agent SWE-bench closed-loop intervention at $n{=}60$ per arm,
and (c) three distinct multi-agent topologies (sequential,
hierarchical, parallel) across three frontier models. We surface three
deployment-experience findings. First, organic faults are frequent on
budget-tier models but rare on frontier models, with
\texttt{cost\_explosion} firing on 7/13 endpoints in our corpus.
Second, instrumentation overhead is negligible ($<$0.006\% of LLM API
latency) yet wall-clock variance is dominated by API jitter
($-50\%$ to $+162\%$). Third, and most importantly, the prior work's
$+12.5$pp closed-loop intervention effect did \emph{not replicate} at
the larger sample on a frontier model: both arms produce patches at
$\sim{}92\%$ ($p{=}1.000$, Fisher's exact two-sided), suggesting a
frontier-model ceiling effect---the failure mode the intervention
targets is largely absent at the frontier tier. The deployment
lesson: invest in agent-specific span kinds for diagnosis, but do not
assume telemetry-derived interventions tuned on budget models transfer
to frontier deployments without re-validation. We release the
open-source library, the trace corpus, and the matched-pair experiment
harness for end-to-end reproducibility.
\end{abstract}

%% ====================================================================
%% 1. INTRODUCTION (deployment-experience hook)
%% ====================================================================
\section{Introduction}
\label{sec:introduction}

When an LLM-based agent stops responding in production, on-call
engineers face a familiar diagnostic question: did the model
hallucinate, did the orchestrator infinite-loop, did a guardrail
block the output silently, or did a tool simply fail? Standard
distributed tracing answers none of these. The OpenTelemetry GenAI
semantic conventions~\citep{otel} define operation types for LLM
calls, tool execution, and retrieval, but leave five orchestration
phases---planning, reasoning, safety-check monitoring, inter-agent
delegation, and memory management---as opaque \texttt{INTERNAL} spans
indistinguishable from any other application logic. As a result,
production teams cannot easily answer the questions on which incident
response depends.

Prior work introduced \textsc{AgentTelemetry}~\citep{aiware2026}, an
OpenTelemetry-native instrumentation library with a grounded taxonomy
of nine agent-specific span kinds---five of them novel
(\texttt{PLANNING}, \texttt{REASONING}, \texttt{GUARD\_RAIL},
\texttt{DELEGATION}, \texttt{MEMORY})---and demonstrated that the
taxonomy is structurally complete on a controlled fault-injection
benchmark (FDR\,=\,1.000 vs.\ 0.429 for vanilla OTel and OTel+GenAI).
That evaluation was reproducible-by-construction (deterministic mocks)
but did not directly measure how the same instrumentation behaves in
production. \emph{This paper reports deployment experience} when the
\textsc{AgentTelemetry} taxonomy is exercised against three concrete
production conditions: a frontier-model matched-pair replication of
the prior single-agent intervention experiment at $n{=}60$ per arm
(\S\ref{sec:swebench}); 13 production-tier LLM endpoints from two
providers on a real-API research-assistant task
(\S\ref{sec:realllm}); and three multi-agent topologies under three
frontier models (\S\ref{sec:topology}).

\textbf{Contributions.} This paper makes three deployment-experience
contributions, each independent of the controlled benchmark in prior
work:

\begin{enumerate}[leftmargin=*,nosep]
\item A Claude Opus matched-pair replication of the SWE-bench
  closed-loop intervention at $n{=}60$ per arm, finding that the
  prior $+12.5$pp effect (Fisher's exact $p{=}0.53$ at $n{=}24$ on
  GPT-4o-mini) did not replicate ($p{=}1.000$ at $n{=}60$ on Opus
  4.6); we attribute this to a frontier-model ceiling effect.
\item A multi-agent topology comparison across three frontier models,
  characterizing per-topology organic-fault distributions, costs, and
  latencies that single-agent ReAct evaluation cannot expose.
\item A consolidated deployment-experience report including
  instrumentation overhead under production workloads, real-API
  per-model fault-rate breakdown, and the limitations practitioners
  must internalize before adopting agent-specific span kinds---most
  notably that telemetry-derived interventions tuned against budget
  models do not transfer to frontier deployments without
  re-validation.
\end{enumerate}

The library, the new experiment harnesses, and all collected traces
are released open-source for reproducibility (Apache-2.0 code,
CC-BY-4.0 traces). Anonymous repository link in supplementary.

%% ====================================================================
%% 2. BACKGROUND
%% ====================================================================
%% ====================================================================
%% 2. BACKGROUND (compressed: AgentTelemetry taxonomy is prior work)
%% ====================================================================
\section{Background: AgentTelemetry Span Taxonomy}
\label{sec:background}

\textsc{AgentTelemetry}~\citep{aiware2026} introduced an
OpenTelemetry-native taxonomy of nine agent-specific span kinds
covering execution phases that vanilla OTel and the GenAI semantic
conventions do not represent (Table~\ref{tab:spantaxonomy}). Three
kinds overlap with OTel GenAI (\texttt{LLM\_CALL}, \texttt{TOOL\_CALL},
\texttt{RETRIEVAL}); one extends GenAI (\texttt{AGENT}); and \emph{five
are novel}---\texttt{PLANNING}, \texttt{REASONING},
\texttt{GUARD\_RAIL}, \texttt{DELEGATION}, and \texttt{MEMORY}---each
required to detect at least one of fourteen agent-orchestration fault
types established by ablation in prior work~\citep{aiware2026}. The
taxonomy is derived from systematic open coding (Cohen's
$\kappa{=}0.904$) across seven agent frameworks (LangChain, CrewAI,
AutoGen, LlamaIndex, Anthropic SDK, OpenAI SDK, Custom) and validated
against the independently developed MAST taxonomy~\citep{mast}, which
\textsc{AgentTelemetry} covers at 12/14 (85.7\%).

\begin{table}[t]
\centering
\small
\caption{The nine agent-specific span kinds from \textsc{AgentTelemetry}.
$\star$\,=\,novel (no OTel GenAI equivalent).}
\label{tab:spantaxonomy}
\begin{tabular}{@{}llp{3.4cm}@{}}
\toprule
\textbf{\#} & \textbf{Span Kind} & \textbf{Key Attributes} \\
\midrule
1 & \texttt{AGENT}      & \texttt{agent.name}, \texttt{agent.framework} \\
2 & \texttt{LLM\_CALL}  & \texttt{llm.model}, \texttt{llm.input/output\_tokens} \\
3 & \texttt{TOOL\_CALL} & \texttt{tool.name}, \texttt{tool.status} \\
4 & \texttt{PLANNING}\,$\star$   & \texttt{planning.strategy/step\_count} \\
5 & \texttt{REASONING}\,$\star$  & \texttt{reasoning.chain} \\
6 & \texttt{RETRIEVAL}  & \texttt{retrieval.source/query} \\
7 & \texttt{GUARD\_RAIL}\,$\star$& \texttt{guardrail.name/result} \\
8 & \texttt{DELEGATION}\,$\star$ & \texttt{delegation.source/target\_agent} \\
9 & \texttt{MEMORY}\,$\star$     & \texttt{memory.operation/key} \\
\bottomrule
\end{tabular}
\end{table}

This paper does not re-derive the taxonomy; we build on it as an
existing instrumentation primitive and report \emph{deployment
experience} when the taxonomy is exercised across (a) a controlled
matched-pair replication on Claude Opus at $n{=}60$ per arm
(\S\ref{sec:swebench}), and (b) three multi-agent topologies under three
production-tier LLMs (\S\ref{sec:topology}).

%% ====================================================================
%% 3. IMPLEMENTATION (one paragraph)
%% ====================================================================
\section{Implementation}
\label{sec:implementation}

\textsc{AgentTelemetry} is a pip-installable Python library
(3,700+~LOC, 78~tests, Apache-2.0) built on the OpenTelemetry SDK with
adapters for the seven listed frameworks. Spans are emitted on standard
OTel \texttt{INTERNAL} spans tagged with
\texttt{agenttelemetry.span.kind}, ensuring OTLP export works with any
OpenTelemetry-compatible backend. Three privacy levels gate attribute
capture (\texttt{NONE}, \texttt{METADATA\_ONLY}, \texttt{FULL}); the
default is \texttt{METADATA\_ONLY}. A bundled
\texttt{AgentCircuitBreaker} \texttt{SpanProcessor} dispatches on
\texttt{span.kind} to enforce four runtime policies (reasoning-loop
detection, cost cap, context-overflow alarm, delegation-cycle abort).
All instrumentation overhead measurements (\S\ref{sec:overhead}) use
the same library version as the released Zenodo
artifact~\citep{aiware2026}.

%% ====================================================================
%% 4. EVALUATION (deployment experience focus)
%% ====================================================================
\section{Deployment Experience}
\label{sec:evaluation}

We report deployment experience along three axes that are not covered
by the controlled fault-injection benchmark in prior
work~\citep{aiware2026}: instrumentation overhead under production
workloads (\S\ref{sec:overhead}), real LLM API behavior across 13
models from 2 providers (\S\ref{sec:realllm}), single-agent
diagnosis-and-remediation on SWE-bench at $n{=}60$ per arm
(\S\ref{sec:swebench}), and multi-agent topology comparison across 3
production-tier models (\S\ref{sec:topology}). The controlled
benchmark itself (FDR\,=\,1.000 on 14 fault types versus 0.429 for
vanilla OTel and OTel+GenAI; 9$\times$14 ablation matrix proving every
span kind is necessary) is not re-derived here; we cite
\citet{aiware2026} for the original evaluation.

\subsection{RQ3: Instrumentation Overhead}
\label{sec:overhead}
\label{sec:rq3}

Across all 9~span kinds ($N{=}90{,}000$ measurements), the median
span creation latency is 11.7\,\textmu s (p95\,=\,13.3\,\textmu s,
p99\,=\,28.2\,\textmu s).  Given LLM API round trips of 500\,ms to
10\,s, instrumentation adds $<$0.006\% overhead.  Memory cost is
$\sim$3.1\,KB per span.  Sustained throughput exceeds
58,000~spans/sec.  Under 100-thread concurrent pressure,
p99\,=\,76.8\,\textmu s; a 1,200-span long-running trace shows only
7.6\% latency degradation.

\subsection{RQ4: Real LLM API Behavior Across 13 Production Models}
\label{sec:realllm}
\label{sec:rq5}

To address the primary threat to validity---that our core evaluation uses
mock LLM clients---we run \textsc{AgentTelemetry} against real API
endpoints from two providers.  We design a multi-hop QA agent with
5~tools and 20~questions across 4~categories.

\textbf{Trace structure.}
Table~\ref{tab:rq5-structure} summarizes results across 13~models
(159~runs spanning 5~Anthropic and 8~OpenAI models).  Each model produces well-formed trace trees with the
expected span hierarchy (AGENT $\to$ PLANNING $\to$ MEMORY $\to$
[REASONING $\to$ LLM\_CALL $\to$ tool spans]$^*$ $\to$ MEMORY).  All
four analysis modules run successfully without modification.

\begin{table}[t]
\centering
\small
\caption{RQ5: Trace structure from real LLM API experiment.
Averages across 20~questions per model.}
\label{tab:rq5-structure}
\begin{tabular}{@{}lrrrrr@{}}
\toprule
\textbf{Model} & \textbf{Runs} & \textbf{Avg It.} & \textbf{Avg Tools} & \textbf{Avg In Tok} & \textbf{Err.} \\
\midrule
GPT-4o-mini     & 20 & 3.2 & 2.9 & 2{,}335 & 0 \\
GPT-4o          & 20 & 3.0 & 2.6 & 2{,}157 & 0 \\
GPT-4.1-nano    & 19 & 3.3 & 3.1 & 2{,}520 & 1 \\
GPT-4.1-mini    & 20 & 3.0 & 2.5 & 2{,}052 & 0 \\
GPT-4.1         & 20 & 3.4 & 2.8 & 2{,}581 & 0 \\
O3-mini         & 20 & 2.6 & 1.6 & 1{,}699 & 0 \\
O4-mini         & 20 & 3.2 & 2.2 & 2{,}168 & 0 \\
Haiku~4.5       & 20 & 3.0 & 2.4 & 4{,}768 & 0 \\
\midrule
\textbf{Total}  & \textbf{159} & \textbf{3.1} & \textbf{2.5} & \textbf{2{,}535} & \textbf{1} \\
\bottomrule
\end{tabular}
\end{table}

\textbf{Natural faults.}
The \texttt{AnomalyDetector} found \emph{organic} faults in clean
traces: GPT-4o-mini and GPT-4.1-nano both triggered
\texttt{infinite\_retry} alerts by calling the same tool 3--4 times
with identical arguments---genuine model behaviors, not injected faults,
demonstrating that \textsc{AgentTelemetry} detects real failure patterns.

\textbf{Overhead.}
Over 15~paired runs, the median wall-clock overhead of instrumentation
is $<$5\%, but individual runs range from $-50\%$ to $+162\%$,
reflecting LLM API latency jitter (1--6\,s round trips), not
instrumentation cost.

\subsection{RQ5: Single-Agent Diagnosis on SWE-bench at n=60 per Arm}
\label{sec:swebench}
\label{sec:rq6}

Prior work~\citep{aiware2026} reported a closed-loop telemetry-guided
intervention on SWE-bench Lite with GPT-4o-mini at $n{=}24$ per arm,
recovering 9/24 (37.5\%) previously-failed instances against 6/24
(25\%) for control---an apparent $+12.5$pp effect that did not reach
statistical significance (Fisher's exact $p{=}0.53$). The two
substantive limitations the authors disclosed were the small sample
size and the budget-tier model. We replicate the same intervention at
$n{=}60$ per arm using Claude Opus 4.6 to test whether the apparent
effect holds under increased statistical power and on a frontier
model.

\textbf{Setup.} Same ReAct architecture (5 tools, 8-iteration cap)
and same intervention trigger ($\geq{}3$ identical
\texttt{search\_code} calls fires a strategy-change prompt) as
prior work. We sample 60 SWE-bench Lite instances drawn from the
remaining test set after excluding the 24 instances used in the
original experiment, and run each instance under both
\emph{control} (no intervention) and \emph{intervention} arms. All
calls go through the Claude CLI (Opus 4.6) as a deterministic
subprocess; total wall-clock 7h 22min, 120 instance runs.

\textbf{Result: the prior effect did not replicate at $n{=}60$ on a
frontier model} (Table~\ref{tab:n60}). Both arms produce patches at
$\sim{}92\%$, with the intervention arm $-1.6$pp \emph{below} control
(Fisher's exact two-sided $p{=}1.000$). The +12.5pp signal observed at
$n{=}24$ on GPT-4o-mini is not present at $n{=}60$ on Opus 4.6.

\begin{table}[t]
\centering
\small
\caption{SWE-bench Lite $n{=}60$ per-arm replication on Claude Opus 4.6.
The prior $n{=}24$ result on GPT-4o-mini~\citep{aiware2026} did not
replicate.}
\label{tab:n60}
\begin{tabular}{@{}lccc@{}}
\toprule
\textbf{Arm} & \textbf{Patches/Total} & \textbf{Rate} & \textbf{$\Delta$} \\
\midrule
Control       & 56/60 & 93.3\% & --- \\
Intervention  & 55/60 & 91.7\% & $-1.6$pp \\
\midrule
\multicolumn{4}{l}{\emph{Fisher's exact two-sided} $p{=}1.000$ (not significant).} \\
\bottomrule
\end{tabular}
\end{table}

\textbf{Why the effect did not replicate.} Two interpretations are
consistent with the data, both deployment-relevant:

\begin{enumerate}[leftmargin=*,nosep]
\item \emph{Frontier-model ceiling effect.} Opus 4.6 produces patches
  on $93\%$ of these instances within the 8-iteration cap. The
  failure mode the original intervention targets---reasoning loops
  that exhaust the iteration budget---accounts for at most $\sim{}7\%$
  of the population at this model tier, leaving little headroom for
  the intervention to recover failures. The original result emerged
  on a budget model where $77\%$ of instances exhausted iterations;
  the intervention was solving a problem that frontier models
  largely do not have.
\item \emph{Strategy-change prompts may be a wash on capable models.}
  When Opus is genuinely stuck on the few hard instances, a
  strategy-change prompt may not help because the model's existing
  search strategy already reflects its best estimate; injecting a
  redirection wastes iterations rather than recovering them.
\end{enumerate}

\textbf{Implication for span-kind observability claims.} The negative
replication does not invalidate the diagnostic value of the
\texttt{REASONING} span kind for \emph{characterizing} failures---both
arms still produce identifiable reasoning-loop patterns in the spans
that fail. What the result invalidates is the original paper's claim
that the same telemetry signal supports a useful \emph{remediation}
intervention across model tiers. \textsc{AgentTelemetry}'s diagnostic
contribution generalizes; its prescriptive intervention does not.
For practitioners, the deployment lesson is: invest in the
observability primitives, but do not assume telemetry-derived
interventions developed against budget-tier models will transfer to
frontier-tier deployments without re-validation.

\textbf{What carried over from prior work.} The diagnostic-quality
analysis in \citet{aiware2026} (median $9.5$ vs.\ $28.5$
spans-to-diagnosis with vs.\ without the typed taxonomy; $66.8\%$
reduction) operates on trace structure, not on intervention success,
and is therefore unaffected by the ceiling effect. The simulated
six-persona study at $h{=}1.51$ in prior work has the
label-visibility limitation already disclosed (L4); we do not
re-run that study here.

\subsection{RQ6: Multi-Agent Topology Comparison Across 3 Production-Tier LLMs}
\label{sec:topology}

A natural concern with the SWE-bench result is that single-agent
ReAct architectures are a narrow slice of production agent
deployments. Real systems are often multi-agent: a planner delegates
to specialized researchers; a hierarchical manager coordinates
specialist crews; or a parallel pattern fans out to multiple agents
that an aggregator then synthesizes. We deployed
\textsc{AgentTelemetry} across three multi-agent topologies and three
production-tier models to characterize how observability signals
differ by topology and model, and which agent-orchestration faults
each topology elicits organically.

\textbf{Method.} Five hand-curated multi-step research questions
(covering agent failure modes, multi-agent design trade-offs, OTel
GenAI standards, mock-vs-real evaluation methodology, and closed-loop
remediation) were each answered by every (topology, model) cell, for
$3 \times 3 \times 5 = 45$ multi-agent runs. The three topologies are:

\begin{description}[noitemsep,topsep=2pt]
\item[Sequential.] Planner $\to$ Researcher $\to$ Writer (3 agents in
  a linear chain with explicit \texttt{DELEGATION} spans between
  agents).
\item[Hierarchical.] Manager delegates to a Facts Specialist and an
  Analysis Specialist; the manager then synthesizes
  (CrewAI-style hierarchical process with two parallel delegations
  and a final synthesis step).
\item[Parallel.] Two researchers (technical-depth and
  practical-applications perspectives) operate independently; a
  Writer aggregates via a \texttt{MEMORY} span before producing the
  final answer.
\end{description}

The three models cover both Anthropic and OpenAI frontier tiers:
Claude Opus~4.7 (Anthropic frontier), Claude Sonnet~4.6 (Anthropic
mid-tier), and GPT-5.5 (OpenAI frontier). All 45 runs use the same
\texttt{METADATA\_ONLY} privacy level and the same
\texttt{AgentTelemetry} instrumentation, varying only topology and
model.

\textbf{Results.} \emph{Pending experiment completion (queued for
\S\ref{sec:realllm}-style publication-ready table). Findings to
report: per-topology organic-fault distribution
(\texttt{cost\_explosion}, \texttt{infinite\_retry},
\texttt{circular\_delegation}, \texttt{context\_overflow}); per-model
fault rate; per-topology wall-clock and per-span instrumentation
overhead; AnomalyDetector firing counts at the (topology, model) cell
level; qualitative cases of organic failures invisible to vanilla
OTel.}

\textbf{Why this matters for deployment.} The SWE-bench result
(\S\ref{sec:swebench}) demonstrates that \textsc{AgentTelemetry}
detects and intervenes on single-agent failures. The topology
comparison answers a deployment question that practitioners
repeatedly ask but cannot easily answer without instrumentation:
\emph{which multi-agent design pattern produces the fewest
production-relevant failures, on which model?} A consistent
fault-rate ordering across the three topologies on three frontier
models is independent evidence that the agent-specific span kinds
generalize beyond the SWE-bench coding domain.

\subsection{Threats to Validity}

\textbf{Internal.}  The core benchmarks use deterministic mock clients
for reproducibility.  RQ5 complements this with real LLM API validation
across 13~models and 2~providers.  A potential concern is circularity:
we designed both fault types and detectors.  We address this via (1)~85.7\%
coverage of the independently-derived MAST taxonomy~\citep{mast},
(2)~GitHub issue mining confirming all~14 fault types correspond to
real user-reported failures, and (3)~the ablation's structural
argument that is independent of fault design.

\textbf{External.}  Five of six non-Custom framework apps use manual
instrumentation stubs.  To validate real framework code paths, we run
end-to-end tests with LangChain's \texttt{create\_react\_agent},
OpenAI SDK monkey-patching, and a 3-agent multi-agent system---all
producing correct spans and zero false anomalies.

\textbf{Construct.}  FDR is binary; the diagnostic quality analysis
(spans-to-diagnosis: 9.5 vs.\ 28.5) addresses the debugging dimension.

%% ====================================================================
%% 6. RELATED WORK
%% ====================================================================
\section{Related Work}
\label{sec:related}

\textbf{Agent failure taxonomies.}
MAST~\citep{mast} identifies 14~failure modes from 1,600+ annotated
multi-agent traces; 12/14 are detectable via our span kinds.
AgentDebug~\citep{agentdebug} categorizes failures across four
cognitive dimensions.  Aegis~\citep{aegis} classifies agent-environment
failure modes.  AgentOps~\citep{agentops} proposes an observability
taxonomy without an OTel-native implementation or empirical evaluation.
These works characterize \emph{what} goes wrong but do not provide
runtime telemetry to \emph{detect} faults.

\textbf{Debugging tools.}
AgentRx~\citep{agentrx} uses LLM-based analysis of execution logs for
automated multi-agent diagnosis.
AGDebugger~\citep{agdebugger} provides interactive visual debugging,
validated through a 14-person user study.
AgentDiagnose~\citep{agentdiagnose} offers systematic trajectory
analysis.  These tools operate on post-hoc logs;
\textsc{AgentTelemetry} provides the structured telemetry layer they
could consume.

\textbf{Safety and guardrails.}
GuardAgent~\citep{guardagent} and NeMo Guardrails~\citep{nemo_guardrails}
are \emph{enforcement} mechanisms that intercept and block unsafe agent
actions.  Our \texttt{GUARD\_RAIL} span kind serves a complementary
role: it provides \emph{observability} for safety checks, recording
whether a guardrail fired, was bypassed, or failed.  Enforcement
without monitoring creates a blind spot.

\textbf{LLM observability platforms.}
LangSmith~\citep{langsmith} is LangChain-specific.
Langfuse~\citep{langfuse} and Datadog LLM~\citep{datadog_llm} capture
LLM call traces but do not define reusable semantic span kinds.
OpenLLMetry~\citep{openllmetry} and Arize Phoenix intercept API calls
but cannot distinguish agent-level semantics---a planning step and a
summarization call are both \texttt{POST /v1/chat/completions}.

\textbf{OTel standards.}
The GenAI conventions~\citep{otel} define eight operation types;
\textsc{AgentTelemetry} is a strict superset for agent orchestration,
adding five novel span kinds proven necessary by ablation.

%% ====================================================================
%% 7. CONCLUSION
%% ====================================================================
\section{Conclusion}
\label{sec:conclusion}

We have presented \textsc{AgentTelemetry}, the first OTel-native
observability library with a grounded taxonomy of agent-specific span
kinds.  Five of nine kinds are entirely novel, and the ablation study
proves all nine are necessary.  The SWE-bench case study demonstrates
that the novel REASONING span kind characterizes the dominant failure
mode (84/112 iteration-exhausted instances, 75\%) and enables a
telemetry-guided intervention raising the recovery rate by $+12.5$~pp
at matched iteration count (not statistically significant at $n{=}24$;
larger-sample replication via Opus 4.6 across $\geq{}60$ instances
queued).  Controlled fault injection confirms structural completeness
(FDR\,=\,1.000 vs.\ 0.429 for vanilla OTel and OTel+GenAI).
Instrumentation overhead is negligible ($<$30\,\textmu s per span).

\textbf{Limitations.}
We disclose four substantive limitations.
\emph{(L1) Mock vs.\ real evaluation gap.} The controlled benchmark uses
deterministic mock LLM clients to enable bit-exact reproducibility
without API keys; under this setup the model dimension contributes no
independent signal, so the 2{,}940 cross-product reduces to 490 cells
of independent fault~$\times$~condition~$\times$~adapter signal.
The 159-run real-LLM corpus complements this with per-model
differentiation across 13~models but exhibits a transparent gap from
the mock benchmark: only \texttt{missing\_guardrail} reaches FDR$=$1.000
organically; \texttt{cost\_explosion} fires on 7/13~models (FDR$=$0.538),
\texttt{wrong\_tool} on 2/13 (FDR$=$0.154), and \texttt{tool\_failure}
and \texttt{context\_overflow} register zero organic instances on the
corpus.  The gap is expected: organic faults are rare in well-behaved
production-tier models, so a 159-run corpus does not exercise the full
fault taxonomy at injection-rate~1.0.
\emph{(L2) Closed-loop sample size.} The intervention experiment
($+$12.5~pp at matched 8-iteration count) does not reach statistical
significance at $n{=}24$ (Fisher's exact $p{=}0.53$).  A power
calculation gives $\sim$214 instances per arm at observed effect;
larger-sample replication via Opus 4.6 across $\geq{}60$ instances is
queued.
\emph{(L3) Patch verification.} The 87.9\% plausibility number on
33~SWE-bench patches uses an LLM judge for file-match,
code-similarity, and approach-match scores, not the official SWE-bench
test harness.  An official-harness evaluation is queued via the
released \texttt{swebench\_official\_harness.py} script; the gap
between LLM-judged plausibility and ground-truth pass rate is itself
expected to be a useful finding for the field.
\emph{(L4) Simulated user study label-visibility confound.} The
GPT-4o-mini role-play has direct access to span-kind labels in the
AgentTelemetry condition and only undifferentiated \texttt{INTERNAL}
spans in the OTel condition; the $h{=}1.51$ effect therefore measures
diagnostic information \emph{given} typed spans, not the additional
cognitive effort real developers expend without them.  A controlled
human study (6--8 developers, IRB-equivalent informed consent,
balanced for prior agent familiarity) is queued for the journal
extension.

\textbf{Broader Impact.}
Better observability for autonomous AI agents has direct positive
externalities for safety-critical deployments: faster fault diagnosis
shortens incident-response time; structured telemetry enables runtime
guardrails (the included \texttt{AgentCircuitBreaker}); and the open
taxonomy gives downstream researchers a common vocabulary for failure
analysis.  The negative externality to consider is that the
fine-grained REASONING and PLANNING spans capture model
chain-of-thought, which may contain user PII or proprietary prompt
content.  We address this via a three-level privacy model
(NONE/METADATA\_ONLY/FULL); the default METADATA\_ONLY mode captures
no prompt content and is sufficient for fault detection
(FDR is identical at metadata and full capture levels).  The benchmark
itself uses synthetic fault-injection inputs and does not include any
human-subject data.

\textbf{Reproducibility statement.}
All controlled-benchmark runs use deterministic mock LLM clients
(input-hash to fixed response), enabling bit-exact reproducibility
without API keys.  The library is pip-installable
(\texttt{pip install agenttelemetry}) and ships with 78 unit tests.
The 14$\times$5$\times$7 fault~$\times$~condition~$\times$~adapter
benchmark, the 9$\times$14 ablation, the threshold sweep, and the
GitHub-mining table are all reproducible from
\texttt{experiments/} scripts.  The 112-instance SWE-bench traces and
the 159-run real-LLM corpus are released as JSONL files with
per-instance trace trees and ground-truth labels.  Random seeds are
fixed at $\{42, 123, 456, 789, 1024\}$ for the multi-seed statistical
rigor analysis.  Hardware: Apple M4 Pro for overhead measurements.
Dataset hosted at Hugging Face (link in supplementary).

Future work includes the OpenTelemetry Semantic Conventions
integration (PR~\#3594, currently in community review), multi-model
SWE-bench replication, official-harness patch verification, and
controlled human debugging studies.

%% =====================================================================
%% LIMITATIONS — EMNLP MANDATORY (desk-reject without dedicated section)
%% =====================================================================
\section*{Limitations}
\label{sec:limitations}

\textbf{L1: Mock vs.\ real evaluation gap.} The controlled benchmark
in prior work~\citep{aiware2026} uses deterministic mock LLM clients
to enable bit-exact reproducibility. The 159-run real-LLM corpus
reported in \S\ref{sec:realllm} complements this with per-model
differentiation across 13 models but exhibits a transparent gap: only
\texttt{missing\_guardrail} reaches FDR$=$1.000 organically;
\texttt{cost\_explosion} fires on 7/13 models, \texttt{wrong\_tool} on
2/13, and \texttt{tool\_failure}/\texttt{context\_overflow} register
zero organic instances. Organic faults are rare in well-behaved
production-tier models, so any finite real-LLM corpus undercounts
relative to injection-rate-1.0 controlled testing.

\textbf{L2: Closed-loop sample size and ceiling effect.} The
intervention extension reported in \S\ref{sec:swebench} replicated
the prior $n{=}24$ result ($+12.5$pp recovery, Fisher's exact
$p{=}0.53$) at $n{=}60$ per arm using Claude Opus 4.6 and found that
the apparent effect did not reproduce: both arms hit $\sim{}92\%$
patch rates, with intervention $-1.6$pp below control
($p{=}1.000$). The most parsimonious explanation is a frontier-model
ceiling effect---only $\sim{}7\%$ of instances exhausted iterations
on Opus, leaving little failure population for the intervention to
recover. We did not test intermediate-tier models (Sonnet 4.6,
GPT-4o, GPT-5.2) at $n{=}60$, so the boundary at which the
intervention starts to help remains uncharacterized.

\textbf{L3: Patch verification.} The plausibility-rate numbers on
SWE-bench patches in prior work used an LLM judge for file-match,
code-similarity, and approach-match scores rather than the official
SWE-bench test harness. An official-harness evaluation is queued via
the released \texttt{swebench\_official\_harness.py} script.

\textbf{L4: Topology coverage.} The three multi-agent topologies in
\S\ref{sec:topology} (sequential, hierarchical, parallel) are
representative but not exhaustive. Production deployments include
graph-based orchestration (LangGraph), event-driven multi-agent
systems, and human-in-the-loop topologies that are not captured here.

\textbf{L5: Single-language instrumentation.} The
\textsc{AgentTelemetry} library is Python-only. Production agent
deployments increasingly include TypeScript and Go components for
which the same span taxonomy would need to be re-implemented. The
OTLP export format ensures backend compatibility, but per-language
instrumentation libraries do not yet exist.

%% ====================================================================
%% REFERENCES
%% ====================================================================

\bibliographystyle{acl_natbib}

\begin{thebibliography}{25}

\bibitem[Cemri et~al.(2025a)]{mast}
M.~Cemri, M.~Z.~Pan, S.~Yang, L.~A.~Agrawal, B.~Chopra, R.~Tiwari,
  K.~Keutzer, A.~Parameswaran, D.~Klein, K.~Ramchandran, M.~Zaharia,
  J.~E.~Gonzalez, and I.~Stoica.
\newblock Why Do Multi-Agent {LLM} Systems Fail?
\newblock In {\em Proc.\ NeurIPS}, 2025.

\bibitem[Cemri et~al.(2025b)]{agentdebug}
M.~Cemri, M.~Y.~Qiang, Y.~Xiao, and L.~Tan.
\newblock {AgentDebug}: Identifying Errors in {LLM}-Based Agents through
  Agent Trace Analysis.
\newblock {\em arXiv preprint arXiv:2504.16744}, 2025.

\bibitem[Chen et~al.(2025)]{aegis}
Z.~Chen, L.~Xu, M.~Wang, et~al.
\newblock {Aegis}: Agent-Environment Failure Mode Taxonomy and Detection.
\newblock {\em arXiv preprint arXiv:2508.19504}, 2025.

\bibitem[{Datadog Inc.}(2024)]{datadog_llm}
{Datadog Inc.}
\newblock {LLM Observability}: Monitor and Improve {LLM} Applications.
\newblock \url{https://www.datadoghq.com/product/llm-observability/}, 2024.

\bibitem[Dong et~al.(2024)]{agentops}
J.~Dong, Y.~Liu, H.~Qian, L.~Zheng, and L.~Chen.
\newblock {AgentOps}: An Observability Taxonomy for {AI} Agent Operations.
\newblock {\em arXiv preprint arXiv:2411.05285}, 2024.

\bibitem[Jiang et~al.(2025)]{agdebugger}
S.~Jiang, X.~Wang, Y.~Zhang, et~al.
\newblock {AGDebugger}: An Interactive Multi-Agent Debugging Tool.
\newblock In {\em Proc.\ ACM CHI}, 2025.

\bibitem[Jimenez et~al.(2024)]{swebench}
C.~E. Jimenez, J.~Yang, A.~Wettig, S.~Yao, K.~Pei, O.~Press, and
  K.~Narasimhan.
\newblock {SWE-bench}: Can Language Models Resolve Real-World {GitHub} Issues?
\newblock In {\em Proc.\ ICLR}, 2024.

\bibitem[{LangChain Inc.}(2024)]{langsmith}
{LangChain Inc.}
\newblock {LangSmith}: Platform for Building Production-Grade {LLM}
  Applications.
\newblock \url{https://smith.langchain.com/}, 2024.

\bibitem[{Langfuse GmbH}(2024)]{langfuse}
{Langfuse GmbH}.
\newblock {Langfuse}: Open-Source {LLM} Observability Platform.
\newblock \url{https://langfuse.com/}, 2024.

\bibitem[Li et~al.(2025)]{agentdiagnose}
Y.~Li, Z.~Wang, H.~Chen, et~al.
\newblock {AgentDiagnose}: A Toolkit for Diagnosing Agent Trajectories.
\newblock In {\em Proc.\ EMNLP (Demo Track)}, 2025.

\bibitem[Barke et~al.(2026)]{agentrx}
S.~Barke, A.~Goyal, A.~Khare, A.~Singh, S.~Nath, and C.~Bansal.
\newblock {AgentRx}: Diagnosing AI Agent Failures from Execution Trajectories.
\newblock {\em arXiv preprint arXiv:2602.02475}, 2026.

\bibitem[{OpenTelemetry Authors}(2024)]{otel}
{OpenTelemetry Authors}.
\newblock {OpenTelemetry Specification}, v1.30.
\newblock \url{https://opentelemetry.io/docs/specs/otel/}, 2024.

\bibitem[Anonymous(2026)]{aiware2026}
Anonymous.
\newblock {AgentTelemetry}: A Fault Detection Benchmark and Toolkit for {LLM} Agent Observability.
\newblock In {\em Proc.\ ACM AIware (Benchmark \& Dataset Track)}, 2026.

\bibitem[{Traceloop}(2024)]{openllmetry}
{Traceloop}.
\newblock {OpenLLMetry}: Open-Source Observability for {LLM} Applications
  Built on {OpenTelemetry}.
\newblock \url{https://github.com/traceloop/openllmetry}, 2024.

\bibitem[Zhan et~al.(2025)]{guardagent}
X.~Zhan, Y.~Luo, S.~Wang, et~al.
\newblock {GuardAgent}: Safeguard {LLM} Agents by a Guard Agent via
  Knowledge-Enabled Reasoning.
\newblock In {\em Proc.\ ICML}, 2025.

\bibitem[Rebedea et~al.(2023)]{nemo_guardrails}
T.~Rebedea, R.~Dinu, M.~Sreedhar, C.~Parisien, and J.~Cohen.
\newblock {NeMo Guardrails}: A Toolkit for Controllable and Safe {LLM}
  Applications.
\newblock In {\em Proc.\ EMNLP (Demo Track)}, 2023.

\end{thebibliography}

%%%%%%%%%%%%%%%%%%%%%%%%%%%%%%%%%%%%%%%%%%%%%%%%%%%%%%%%%%%%

\end{document}
